\begin{theorem}[Riccati lists and solution counts]
\label{thm:riccati-cross-ratio}
Let $D\ge1$, and let $E$ have characteristic zero or $p>D$.
Consider one fixed equation
\begin{equation}
 a(X)P'=b_0(X)+b_1(X)P+b_2(X)P^2,\qquad a\ne0,
 \label{eq:fixed-riccati}
\end{equation}
with coefficients in $E[X]$, with no restriction on their degrees.
On any $n$ distinct evaluation points, the number of distinct solutions
of degree at most $D$ agreeing with an arbitrary received word in at
least $A>D$ coordinates is at most
\begin{equation}
 \left\lfloor\frac{n}{A-D}\right\rfloor.
 \label{eq:riccati-list}
\end{equation}
This bound is attained over prime fields at unbounded lengths and
positive limiting rate and agreement gap. If $b_2\ne0$, the total
number of degree-at-most-$D$ polynomial solutions is at most
\begin{equation}\label{eq:riccati-total-count}
 C_D=\begin{cases}
 D+2,&\text{in characteristic zero or }p>D+1,\\
 2D+2,&\text{if }p=D+1.
 \end{cases}
\end{equation}
The bound $D+2$ is sharp. The boundary bound $2D+2$ is attained
at $p=2,D=1$.
\end{theorem}
\begin{proof}
\emph{Total candidate count.}
Assume $b_2\ne0$. Work at a transcendental Taylor center $t$, where
$a(t)b_2(t)\ne0$, and write the equation as
$P'=F(X,P)$, with $F=\alpha+\beta u+\gamma u^2$ and $\gamma\ne0$.
Put $\Delta=\partial_X+F\partial_u$ and $F_j=\Delta^j u$.
For $j<p$, or in characteristic zero, $F_j$ has $u$-degree $j+1$
and leading coefficient $j!\gamma^j$. This follows inductively:
$F\partial_u$ gives the highest degree and $\partial_X$ gives a
lower degree. In characteristic zero or $p>D+1$, every solution
satisfies
\[
 F_{D+1}(t,P(t))=P^{(D+1)}(t)=0.
\]
This is a nonzero polynomial of degree $D+2$ in $P(t)$.
Distinct polynomial solutions have distinct values at the
transcendental $t$, proving the first count.

If $p=D+1$, use instead the truncated reconstruction
\[
 T(U,X)=\sum_{j=0}^D F_j(t,U)\frac{(X-t)^j}{j!}.
\]
Every candidate is $T(P(t),X)$. Its $U$-degree is $D+1$, with
leading coefficient $\gamma(t)^D(X-t)^D$. The residual
$a\partial_XT-b_0-b_1T-b_2T^2$ therefore has $U$-degree exactly
$2D+2$, with nonzero leading coefficient
$-b_2(X)\gamma(t)^{2D}(X-t)^{2D}$. Any nonzero $X$-coefficient
of this residual is a nonzero polynomial in $U$ of degree at most
$2D+2$ vanishing at every candidate value $P(t)$. This proves the
boundary count. In particular, in positive characteristic the total
count is at most $2D+2\le2p$.

\emph{Lists by intersection multiplicities.}
For $M\ge2$ distinct candidates and a domain coordinate $x$, put
\[
 C_x=\sum_{i<j}\operatorname{ord}_x(P_i-P_j).
\]
The root bound gives $\sum_xC_x\le D\binom M2$. If $h$ candidates
match a given received value at $x$, we claim
\begin{equation}\label{eq:riccati-local-multiplicity}
 (h-1)(M-1)\le2C_x.
\end{equation}
If $b_2=0$, differences solve one homogeneous linear equation.
The ratio of two nonzero differences has derivative zero and
rational degree at most $D<p$, and hence is constant. The same is
true in characteristic zero. Thus the candidates have form $P_0+cW$,
and their values at $x$ are either all distinct or all equal;
\eqref{eq:riccati-local-multiplicity} follows immediately.

Now suppose $b_2\ne0$. Differences satisfy
\[
 a(P_i-P_j)'=(b_1+b_2(P_i+P_j))(P_i-P_j),
\]
so every cross ratio has derivative zero. In characteristic zero
it is constant and different from $0,1,\infty$. A colliding pair
and two candidates outside its value class would force that cross
ratio to be one. Therefore either all values are distinct or one
value occurs at least $M-1$ times. For
$h\in\{0,1,M-1,M\}$, the inequality follows from
$C_x\ge\binom h2$.

In positive characteristic a cross ratio need not be constant.
However, any zero of a nonzero rational function with derivative
zero has order divisible by $p$. Choose a colliding pair and two
candidates outside its value class. The denominator of their
cross-ratio identity is nonzero at $x$, and the numerator of the
cross ratio minus one, up to sign, is the product of the two
within-pair differences. If the two outside values were distinct,
this product would have positive order at most $D<p$, impossible.
Consequently every collision pattern other than those already
handled has exactly two value classes, both of size at least two.

Let the matching class have size $h$ and the other class have size
$b=M-h$. Let $\mu,\nu$ be the minimum positive vanishing orders of
differences within these classes. The same cross-ratio argument
for every choice of two within-class pairs gives $\mu+\nu\ge p$.
Write $A_0=\binom h2$, $B_0=\binom b2$, so
$C_x\ge A_0\mu+B_0\nu$. If $h\le b$, already
\[
 2(A_0+B_0)-(M-1)(h-1)=(b-1)(b-h+1)\ge0.
\]
If $h\ge b$, minimizing with $\mu,\nu\ge1$ and
$\mu+\nu\ge p$ gives $C_x\ge A_0+(p-1)B_0$, and hence
\[
 2C_x-(M-1)(h-1)\ge(b-1)(pb-M+1)\ge0,
\]
because $b\ge2$ and the total-count bound gives $M\le2p$.
An unmatched received value has $h=0$. This proves
\eqref{eq:riccati-local-multiplicity} in all cases.
Summing it gives $AM\le n+DM$, proving
\eqref{eq:riccati-list}. The case $M=1$ is immediate, and the same
proof allows real agreement thresholds.

For sharpness of $D+2$, choose $a$ split and squarefree of degree
$D+1$. The equation $aP'=a'P-P^2$ has the $D+2$ solutions $0$ and
$a/(X-r)$, one for each root $r$ of $a$. At the boundary $p=2,D=1$,
$(X^2+X)P'=P^2+P$ has exactly $0,1,X,X+1$ as its four solutions.
\end{proof}
